Circuit representation models such as DeepGate are trained to capture what each gate computes, so gates computing the same
Boolean function should get the same predicted signal probability (the chance of outputting 1) and nearby embeddings however the
surrounding logic is rewritten. We test this premise with \audit, which rewrites circuits by logic synthesis and by
adding redundant logic and proves with SAT that each rewrite preserves function and which gates correspond. The
premise fails. On \NTestCones\ EPFL test circuits, re-synthesis keeps a DeepGate-style model accurate yet moves its
prediction for a fixed function by about as much as it errs, enough to change the phase hints it gives a SAT
solver, at every size up to \ScMaxParams\ parameters. An XOR-based rewrite unseen in training breaks every model's
predictions, and for five released encoders, two of them trained for equivalence, it at least halves
how often a rewritten gate's nearest original gate in embedding space is its proven equivalent. Training on
certified rewrites steadies predictions, mostly against redundancy, but only aligning proven-equivalent gates
(\ours) raises this rate on seen rewrites from \BaseTopRS\% to \OursTopRS\%, beyond what 4$\times$ more unaligned rewrites
reach; fine-tuning released DeepGate2 with \ours\ reaches \FtEqTopSeen\% (\FtCtrlTopSeen\% with its original loss).
No tested repair survives the unseen rewrite.
